\documentclass[aspectratio=169,11pt]{beamer}
\usetheme{Boadilla}
\usecolortheme{seahorse}
\usepackage{listings}
\usepackage{amsmath,amssymb,booktabs}
\definecolor{codeblue}{RGB}{25,65,125}
\setbeamertemplate{navigation symbols}{}
\setbeamertemplate{footline}[frame number]
\setbeamercolor{structure}{fg=codeblue}
\lstset{language=Python,basicstyle=\ttfamily\small,keywordstyle=\color{codeblue}\bfseries,commentstyle=\color{gray},stringstyle=\color{teal},showstringspaces=false,columns=fullflexible,keepspaces=true,breaklines=true,tabsize=4}
\title{A Closure-Based Lambda Interpreter}
\subtitle{Reading \texttt{lambda1.py}}
\author{CS413}
\date{September 22, 2026}
\begin{document}
\begin{frame}\titlepage\end{frame}

\begin{frame}[fragile]{The evaluator's interface}
\begin{lstlisting}
def d0exp_evaluate\
(dexp: d0exp, denv: d0env = ENVnil()) -> d0val:
    """
    [dexp] may contain free variables
    """
\end{lstlisting}
\bigskip
\begin{itemize}
\item An expression describes a computation.
\item An environment gives values to variable names.
\item Evaluation produces a value, including a closure for a function.
\item The interpreter uses call-by-value and lexical scope.
\end{itemize}
\medskip
The \texttt{type} aliases in this file require Python 3.12 or later.
\end{frame}

\begin{frame}{The expression language}
\begin{center}
\begin{tabular}{ll}\toprule
Constructors & Meaning\\\midrule
\texttt{D0Eint, D0Ebtf} & Integer and Boolean literals\\
\texttt{D0Evar} & Variable lookup\\
\texttt{D0Elam, D0Efix} & Ordinary and recursive functions\\
\texttt{D0Eapp} & Function application\\
\texttt{D0Elet} & Local, nonrecursive binding\\
\texttt{D0Eop1, D0Eop2} & Integer operations and comparisons\\
\texttt{D0Eif0} & Boolean conditional\\
\texttt{D0Epair, D0Epfst, D0Epsnd} & Pairs and projections\\\bottomrule
\end{tabular}
\end{center}
\medskip
The file consumes abstract syntax trees. It does not parse source text.
\end{frame}

\begin{frame}[fragile]{Constructing an expression}
\[
(\lambda x.\,x+1)\;41
\]
\begin{lstlisting}
increment = D0Elam("x",
    D0Eop1("+1", D0Evar("x")))

term = D0Eapp(increment, D0Eint(41))

d0exp_evaluate(term)  # D0Vint(arg1=42)
\end{lstlisting}
\bigskip
\texttt{D0Eapp} constructs syntax. \texttt{d0exp\_evaluate} executes it.\\[0.6em]
\texttt{@dataclass} supplies constructors, representations, and equality.
\end{frame}

\begin{frame}[fragile]{Expressions and values have different roles}
\begin{columns}[T]
\begin{column}{0.49\textwidth}
\begin{lstlisting}
@dataclass
class D0Elam(D0E000):
    arg1: dvar
    arg2: d0exp
    ctag = "D0Elam"
\end{lstlisting}
\medskip
A lambda expression stores its parameter and body.
\end{column}
\begin{column}{0.49\textwidth}
\begin{lstlisting}
@dataclass
class D0Vlam(D0V000):
    arg1: d0env
    arg2: D0Elam
    ctag = "D0Vlam"
\end{lstlisting}
\medskip
A closure stores an environment together with that expression.
\end{column}
\end{columns}
\bigskip
Other values are \texttt{D0Vint}, \texttt{D0Vbtf}, \texttt{D0Vpair}, and \texttt{D0Vfix}.\\[0.6em]
The evaluator dispatches with \texttt{isinstance}, rather than \texttt{ctag}.
\end{frame}

\begin{frame}[fragile]{Environments are linked bindings}
\begin{lstlisting}
@dataclass(frozen=True)
class ENVcns(ENV000):
    arg1: dvar
    arg2: d0val
    arg3: d0env
    ctag = "ENVcns"

outer = ENVcns("x", D0Vint(10), ENVnil())
inner = ENVcns("x", D0Vint(20), outer)
\end{lstlisting}
\medskip
\texttt{ENVnil()} ends the chain. Each new binding retains the outer chain.\\[0.6em]
The links are frozen. The referenced values and syntax trees are not frozen.
\end{frame}

\begin{frame}[fragile]{Lookup chooses the nearest binding}
\begin{lstlisting}
def d0env_search(denv: d0env, dvar: dvar) -> d0val:
    while True:
        if isinstance(denv, ENVcns):
            if dvar == denv.arg1:
                return denv.arg2
            else:
                denv = denv.arg3; continue
        else:
            return D0V000() # HX: this indicates an error
\end{lstlisting}
\medskip
For the preceding environments, lookup of \texttt{x} gives 20 in \texttt{inner},\\
but still gives 10 in \texttt{outer}. A missing name yields \texttt{D0V000()}.
\end{frame}

\begin{frame}[fragile]{Evaluating a lambda creates a closure}
Evaluator excerpt:
\begin{lstlisting}
elif isinstance(dexp, D0Elam):
    return D0Vlam(denv, dexp)
elif isinstance(dexp, D0Efix):
    return D0Vfix(denv, dexp)
\end{lstlisting}
\bigskip
\begin{itemize}
\item The closure captures the environment at the function's definition.
\item Creating a closure does not evaluate its body.
\item The body can use captured bindings when a later call executes it.
\end{itemize}
\medskip
No substitution or copying of the function body occurs here.
\end{frame}

\begin{frame}[fragile]{Application of an ordinary closure}
Excerpts from \texttt{f0\_D0Eapp}:
\begin{lstlisting}
dfun = d0exp_evaluate(dexp.arg1, denv)
darg = d0exp_evaluate(dexp.arg2, denv)
\end{lstlisting}
For a \texttt{D0Vlam} value:
\begin{lstlisting}
dlam = dfun.arg2
denv_new = ENVcns(dlam.arg1, darg, dfun.arg1)
return d0exp_evaluate(dlam.arg2, denv_new)
\end{lstlisting}
\medskip
The caller's environment evaluates the function and argument, in that order.\\[0.6em]
The \alert{captured environment} evaluates the body, with the parameter added.
\end{frame}

\begin{frame}[fragile]{Lexical scope in a worked example}
Readable notation for a nested AST:
\begin{lstlisting}[language={}]
let x = 10 in
  let f = (lambda y. x + y) in
    let x = 100 in
      f(2)
\end{lstlisting}
\bigskip
\begin{tabular}{ll}
When \texttt{f} is defined & Its closure captures \texttt{x = 10}.\\[0.7em]
When \texttt{f(2)} is called & The caller has \texttt{x = 100}.\\[0.7em]
When the body runs & It finds \texttt{y = 2}, then captured \texttt{x = 10}.\\
\end{tabular}
\par
\bigskip
\alert{Result: \texttt{D0Vint(12)}}
\end{frame}

\begin{frame}[fragile]{A let evaluates its initializer before binding}
Body of \texttt{f0\_D0Elet}:
\begin{lstlisting}
dval = d0exp_evaluate(dexp.arg2, denv)
denv_new = ENVcns(dexp.arg1, dval, denv)
return d0exp_evaluate(dexp.arg3, denv_new)
\end{lstlisting}
\bigskip
In \texttt{D0Elet(name, initializer, body)}:
\begin{itemize}
\item The initializer uses the old environment.
\item Only the body sees the new binding.
\item Even an unused initializer executes.
\end{itemize}
\medskip
Example: with outer \texttt{x = 10}, \texttt{let x = x+1 in x} yields 11.\\
The outer binding remains 10. This form of let is nonrecursive.
\end{frame}

\begin{frame}[fragile]{A recursive function names itself}
\begin{lstlisting}
@dataclass
class D0Efix(D0E000):
    arg1: dvar
    arg2: dvar
    arg3: d0exp
    ctag = "D0Efix"
\end{lstlisting}
\bigskip
\texttt{D0Efix("f", "x", body)} represents $\mathrm{fix}\ f(x).\,body$.\\[0.7em]
\texttt{arg1} names the function itself, \texttt{arg2} names its parameter,\\
and \texttt{arg3} holds its body.\\[0.7em]
Evaluation creates \texttt{D0Vfix(denv, dexp)}. A call adds the self binding.
\end{frame}

\begin{frame}[fragile]{Application of a recursive closure}
The \texttt{D0Vfix} branch of \texttt{f0\_D0Eapp}:
\begin{lstlisting}
dfix = dfun.arg2
denv_new0 = dfun.arg1
denv_new1 = \
    ENVcns(dfix.arg1, dfun, denv_new0)
denv_new2 = \
    ENVcns(dfix.arg2, darg, denv_new1)
return d0exp_evaluate(dfix.arg3, denv_new2)
\end{lstlisting}
\medskip
Lookup order: parameter, function's own name, captured bindings.\\[0.6em]
Binding the function name to \texttt{dfun} lets the body call the same closure.\\[0.6em]
If the parameter and function names coincide, the parameter shadows it.
\end{frame}

\begin{frame}[fragile]{Factorial as an abstract syntax tree}
\begin{lstlisting}
factorial = D0Efix("fact", "n",
    D0Eif0(
        D0Eop2("<=", D0Evar("n"), D0Eint(0)),
        D0Eint(1),
        D0Eop2("*", D0Evar("n"),
            D0Eapp(D0Evar("fact"),
                D0Eop1("-1", D0Evar("n"))))))

d0exp_evaluate(D0Eapp(factorial, D0Eint(4)))
# D0Vint(arg1=24)
\end{lstlisting}
\medskip
Each call adds a new \texttt{n} binding and the same \texttt{fact} closure.\\[0.6em]
The calls descend through 4, 3, 2, 1, 0. Multiplication finishes as they return.
\end{frame}

\begin{frame}[fragile]{Conditionals evaluate only the selected branch}
Body of \texttt{f0\_D0Eif0}:
\begin{lstlisting}
dexp1 = d0exp_evaluate(dexp.arg1, denv)
if not isinstance(dexp1, D0Vbtf):
    raise TypeError(f"D0Vbtf(...) expected: {dexp1}")
if dexp1.arg1:
    return d0exp_evaluate(dexp.arg2, denv)
else:
    return d0exp_evaluate(dexp.arg3, denv)
\end{lstlisting}
\medskip
Despite its name, \texttt{D0Eif0} expects a Boolean, not an integer zero test.\\[0.6em]
In factorial, the base case avoids evaluating the recursive branch.
\end{frame}

\begin{frame}{Primitive operations check value kinds}
\begin{center}
\begin{tabular}{lll}\toprule
Operators & Operand values & Result value\\\midrule
\texttt{+1, -1} & One \texttt{D0Vint} & \texttt{D0Vint}\\
\texttt{+, -, *, /} & Two \texttt{D0Vint}s & \texttt{D0Vint}\\
\texttt{<, >, <=, >=, ==, !=} & Two \texttt{D0Vint}s & \texttt{D0Vbtf}\\\bottomrule
\end{tabular}
\end{center}
\bigskip
\begin{itemize}
\item Binary operations evaluate both operands from left to right.
\item The interpreter's \texttt{/} uses Python's floor division \texttt{//}.\\
For example, $-7/3$ produces \texttt{D0Vint(-3)}.
\item Wrong operand kinds and unsupported operators raise \texttt{TypeError}.
\end{itemize}
\end{frame}

\begin{frame}[fragile]{Pairs evaluate both components}
The \texttt{D0Epair} branch:
\begin{lstlisting}
dval1 = d0exp_evaluate(dexp.arg1, denv)
dval2 = d0exp_evaluate(dexp.arg2, denv)
return D0Vpair(dval1, dval2)
\end{lstlisting}
\medskip
\texttt{D0Epfst(e)} and \texttt{D0Epsnd(e)} evaluate \texttt{e}, require a\\
\texttt{D0Vpair}, then select one component.\\[0.7em]
Pairs can contain mixed kinds, nested pairs, and closures.
\begin{lstlisting}
d0exp_evaluate(D0Epfst(
    D0Epair(D0Eint(7), D0Ebtf(True))))
# D0Vint(arg1=7)
\end{lstlisting}
\end{frame}

\begin{frame}[fragile]{Call-by-value makes unused computations matter}
Let \texttt{bad} be an expression that divides by zero:
\begin{lstlisting}
bad = D0Eop2("/", D0Eint(1), D0Eint(0))

D0Eapp(D0Elam("x", D0Eint(7)), bad)
D0Elet("x", bad, D0Eint(7))
D0Epfst(D0Epair(D0Eint(7), bad))
\end{lstlisting}
Evaluating any of these three expressions raises \texttt{ZeroDivisionError}.
\par
\bigskip
By contrast, this conditional produces \texttt{D0Vint(7)}:
\begin{lstlisting}
D0Eif0(D0Ebtf(True), D0Eint(7), bad)
\end{lstlisting}
\end{frame}

\begin{frame}{Errors and implementation boundaries}
\begin{itemize}
\item An unbound variable returns the sentinel \texttt{D0V000()}.\\
It can propagate as a result or become a pair component.
\item Invalid application, projection, or operand kinds raise \texttt{TypeError}.
\item Application evaluates the argument before checking the function kind.\\
An argument's error can therefore occur first.
\item Dataclass annotations assume correctly formed ASTs and environments.\\
They do not validate constructor arguments.
\item Evaluation uses Python recursion, including for tail calls.\\
Deep calls can exceed the recursion limit. Some terms never terminate.
\end{itemize}
\end{frame}

\begin{frame}[fragile]{Questions for tracing the interpreter}
\begin{enumerate}
\item What happens if application extends the caller's environment instead of
\texttt{dfun.arg1}? Revisit the example with \texttt{x = 10} and \texttt{x = 100}.
\item Why does \texttt{D0Vfix} application add the function's own binding?
\item What does this expression return?
\end{enumerate}
\begin{lstlisting}
D0Eapp(
    D0Eapp(D0Elam("x", D0Elam("y", D0Evar("x"))),
           D0Eint(10)),
    D0Eint(20))
\end{lstlisting}
\medskip
Relevant examples and checks: \texttt{TEST/test03\_lambda1.py}\\
Run with \texttt{python3 -B TEST/test03\_lambda1.py}.
\end{frame}
\end{document}
